\documentclass[8pt, a4paper]{extarticle}
\usepackage[
    a4paper, 
    landscape, 
    margin=0.2in,
    footskip=0.1in,
    headsep=0in,
    headheight=0in
]{geometry}
\usepackage{graphicx}
\usepackage{ctex}
\usepackage{multicol}
\usepackage{amsmath,amssymb,amsfonts,amsthm}
\usepackage{enumitem}
\usepackage{pifont}
\usepackage{siunitx}
\usepackage{xcolor}

\usepackage{tabularx}
\usepackage{booktabs}
\usepackage{multirow}
\renewcommand{\tabularxcolumn}[1]{m{#1}}
\newcolumntype{C}[1]{>{\hsize=#1\hsize\centering\arraybackslash}X}

\usepackage{fix-cm}
\DeclareMathSizes{7}{7}{4}{4}

\setlength{\columnsep}{0.5cm}
\setlength{\parindent}{0pt}
\setlist{noitemsep}
\setlength{\jot}{0.5pt}
\setcounter{secnumdepth}{0}
\pagestyle{empty}

\usepackage[tiny,compact]{titlesec}
\titleformat{\section}{\bfseries\normalsize}{\thesection}{0.5em}{}
\titlespacing*{\section}{0pt}{1ex}{0.5ex}
\titleformat{\subsection}{\bfseries\small}{\thesubsection}{0.4em}{}
\titlespacing*{\subsection}{0pt}{0.5ex}{0.2ex}
\titleformat{\subsubsection}{\normalfont\small\bfseries\itshape}{\thesubsubsection}{0.3em}{}
\titlespacing*{\subsubsection}{0pt}{0.4ex plus .1ex minus .1ex}{0.2ex}
\titleformat{\paragraph}[runin]{\normalfont\small\bfseries}{\theparagraph}{0.5em}{}[.]
\titlespacing*{\paragraph}
  {0pt}
  {0.3ex plus 0.1ex minus 0.1ex}
  {0.5em} 

\newcommand*{\dif}{\mathop{}\!\mathrm{d}}

\DeclareMathOperator{\Res}{Res}
\DeclareMathOperator{\FFT}{FFT}

\makeatletter
\renewcommand{\maketitle}{
  \begin{center}
    \vspace{-2.5em} 
    
    {\Large \bfseries \@title} \\[0.4ex] 
    
    {\scriptsize \@author \quad \text{|} \quad \@date} 
    
    \vspace{0.2ex}
    \hrule 
    
    \vspace{-1.0em} 
  \end{center}
}
\makeatother

\title{奇襲 \textcolor{blue}{D}igital \textcolor{blue}{S}ignal \textcolor{blue}{P}rocessing \ding{43}}
\author{luisleee}
\date{\today}

\begin{document}
\small
\begin{multicols}{4}

\maketitle
\section{基本约定}
记连续时间信号为\(x(t)\)。记离散时间信号即序列为\(x(n)\)，其Z变换为\(X(z)\)，其DTFT为\(X(e^{j\omega})\)，DFT为\(X(k)\)。记周期序列为\(\tilde{x}(n)\)，也可以表示\(x(n)\)的周期延拓，其DFS为\(\tilde{X}(k)\)。所有提到的系统都默认为线性时不变系统。\(\Leftrightarrow\)表示变换对，代表的变换具体视上下文而定。文中的\(\delta\)对离散时间序列表示Kronecker \(\delta\)，对连续时间序列表示Dirac \(\delta\)。复分析相关定理假设曲线是简单的，\(\Res[f(z),z_0]\)表示函数\(f(z)\)在点\(z_0\)处的留数。Z变换默认使用双边变换，性质部分忽略收敛域。对有限长序列进行操作时，当下标超出定义范围，默认按周期延拓处理。用\(\otimes\)表示有限长序列的循环卷积，卷积长度默认相同。\(u(n)\)表示单位阶跃函数，即\(n\geq 0\)时为\(1\)，否则为\(0\)。
\section{基本性质}

\subsection{旋转因子 (单位根)}
记\(N\)次单位根（旋转因子）为 \(W_N = e^{-j2\pi/N}\)

分圆方程 \(\sum_{n=0}^{N-1}W_N^n = 0\)

完备正交性\(\frac{1}{N}\sum_{n=0}^{N-1}W_N^{nk} = \delta(k),\quad k\in \{0,1,\ldots,N-1 \}\)

\subsection{复变函数与留数}
Cauchy积分公式
\( 2\pi j f^{(n)}(z_0) = n! \oint_\Gamma {f(z)(z - z_0)}^{-n-1}\dif z \)

留数定理\(\oint_\Gamma f(z)\dif z = 2\pi j \sum_k \Res[f(z), z_k]\)

极点留数求法：若 \(z_0\) 为 \(f(z)\) 的 \(n\) 阶极点
\[ \Res[f(z), z_0] = \frac{1}{(n-1)!} \lim_{z \to z_0} \frac{\dif^{n-1}}{\dif z^{n-1}} \left[ (z - z_0)^n f(z) \right] \]
特别地，对一阶极点：\(\Res[f(z), z_0] = \lim_{z \to z_0} (z - z_0) f(z)\)。

\subsection{信号抽样}
连续信号\(x_a(t)\)按取样间隔\(T_s\)取样，频率\(f_s=1/T_s\)，得离散信号\(x(n)=x_a(nT_s)\)。取样函数\(p(t)=\sum_{n=-\infty}^{+\infty}\delta(t-nT_s)\)，理想取样信号\(\hat{x}_a(t)=x_a(t)p(t)\)。

模拟抽样角频率\(\Omega_s=2\pi/T_s\)，频谱展开为：
\[ \hat{X}_a(j\Omega)=\Omega_s\sum_{n=-\infty}^{+\infty}X_a(j(\Omega-n\Omega_s)) \]

Nyquist取样定理：取样频率\(f_s \geq 2f_m\) (\(f_m\)是原信号最高频率) 时，可通过理想低通滤波器(理想内插)无失真恢复原信号。

时域频域取样定理：时域取样对应频域周期延拓；频域取样对应时域周期延拓。

模拟角频率 \(\Omega\) (\(\mathrm{rad/s}\)) 与数字角频率 \(\omega\) (\(\mathrm{rad}\)) 关系为 \(\omega = \Omega T_s = \Omega/f_s\)

\section{DFS}
周期为\(N\)的序列\(\tilde{x}(n)\)的DFS变换对
\[
\tilde{X}(n)=\sum_{n=0}^{N-1}\tilde{x}(n)W_N^{nk}
\]
\[
\tilde{x}(n)=\frac{1}{N}\sum_{k=0}^{N-1}\tilde{X}(k)W_N^{-nk}
\]

\section{DTFT}
\subsection{定义}
序列\(x(n)\)的DTFT
\[
X(e^{j\omega}) = \sum_{n=-\infty}^{+\infty}x(n)e^{-jn\omega}
\]
反变换IDTFT
\[
x(n) = \frac{1}{2\pi}\int_{-\infty}^{+\infty}X(e^{j\omega})e^{jn\omega}\dif \omega
\]
\subsection{性质}
线性
\(ax(n)+y(n)\Leftrightarrow aX(e^{j\omega}) + Y(e^{j\omega})\)

共轭
\(x^*(n)\Leftrightarrow X^*(e^{-j\omega})\)

奇偶
\(x(-n)\Leftrightarrow X(e^{-j\omega})\)

周期
\(X(e^{j\omega})=X(e^{j(\omega+2\pi)})\)

时移和频移
\begin{gather*}
x(n-n_0)\Leftrightarrow e^{-jn_0\omega}X(e^{j\omega})\\
e^{jn\omega_0}x(n)\Leftrightarrow X(e^{j(\omega-\omega_0)})
\end{gather*}

频域微分
\[
n x(n)\Leftrightarrow j\frac{\dif X(e^{j\omega})}{\dif \omega}
\]

周期冲激序列
\[
\sum_{k=-\infty}^{+\infty}\delta(n-kN)\Leftrightarrow \frac{2\pi}{N}\sum_{k=-\infty}^{+\infty}\delta\left(\omega-\frac{2\pi k}{N}\right)
\]

卷积定理
\begin{gather*}
x(n)\ast y(n)\Leftrightarrow X(e^{j\omega})Y(e^{j\omega})\\
x(n)y(n)\Leftrightarrow \frac{1}{2\pi} X(e^{j\omega})\ast Y(e^{j(\omega)})
\end{gather*}

内积
\[
\sum_{n=-\infty}^{+\infty}x(n)y^*(n)=\frac{1}{2\pi}\int_{-\pi}^{\pi}X(e^{j\omega})Y^*(e^{j\omega})\dif \omega
\]

Parseval定理
\[
\sum_{n=-\infty}^{+\infty}|x(n)|^2=\frac{1}{2\pi}\int_{-\pi}^{\pi}\left|X(e^{j\omega})\right|^2\dif \omega
\]

\section{Z变换}
\subsection{定义}
序列\(x(n)\)的Z变换
\(
X(z)=\sum_{n=-\infty}^{+\infty}x(n)z^{-n}
\)
\subsection{性质}
将\(z=e^{j\omega}\)代入\(X(z)\)得DTFT

反变换，\(C\)“内部”包含\(X(z)\)所有极点
\[
x(n)=\frac{1}{2\pi j}\oint_{C}X(z)z^{n-1}\dif z=\sum_{p_k}\Res[X(z)z^{n-1}, p_k]
\]

线性
\(
ax(n)+y(n)\Leftrightarrow aX(z)+Y(z)
\)

共轭
\(
x^*(n) \Leftrightarrow X^*(z^*)
\)

时移
\(
x(n-n_0) \Leftrightarrow z^{-n_0}X(z)
\)

尺度变换
\(
a^n x(n) \Leftrightarrow X(z/a)
\)

微分
\(
nx(n)\Leftrightarrow -z\frac{\dif}{\dif z}X(z)
\)

卷积
\(
x(n)\ast y(n)\Leftrightarrow X(z)\cdot Y(z)
\)

\subsection{常见Z变换}
{
\normalsize
\renewcommand{\arraystretch}{1.4}
\begin{tabularx}{\linewidth}{|C{1}|C{1.4}|C{0.6}|}
    \toprule
    序列 & Z变换 & 收敛域 \\
    \midrule
    \(\delta(n)\) & \(1\)  & \(\mathbb{C}\) \\
    \hline
    \(u(n)\) & \(\frac{1}{1-z^{-1}}\) & \(|z|>1\) \\
    \hline
    \(-u(-n-1)\) & \(\frac{1}{1-z^{-1}}\) & \(|z|<1\) \\
    \hline
    \(\sin(\omega_0 n)u(n)\) & \(\frac{\sin(\omega_0)z^{-1}}{1-(2\cos \omega_0)z^{-1}+z^{-2}}\) & \(|z|>1\) \\
    \hline
    \(\cos(\omega_0 n)u(n)\) & \(\frac{1-\cos(\omega_0)z^{-1}}{1-(2\cos \omega_0)z^{-1}+z^{-2}}\) & \(|z|>1\) \\
    \hline
    \(n u(n)\) & \(\frac{z^{-1}}{{(1-z^{-1})}^2}\) & \(|z|>1\) \\
    \hline
    \(-n u(-n-1)\) & \(\frac{z^{-1}}{{(1-z^{-1})}^2}\) & \(|z|<1\) \\
    \bottomrule
\end{tabularx}
}


\section{离散系统性质}
\paragraph{因果性} 充要条件：\(h(n) = 0 \ (n < 0)\)；频域表现：\(H(z)\)收敛域包含\(\infty\)。
\paragraph{稳定性} 充要条件：\(\sum_{n=-\infty}^{+\infty}|h(n)|<\infty\)；频域表现：\(H(z)\)收敛域包含单位圆。
\paragraph{因果稳定系统分解} 可分解为全通与最小相位系统级联：\(H(z) = H_{\text{min}}(z) \cdot H_{\text{ap}}(z)\)。
\paragraph{全通系统} 因果稳定，幅频响应恒为常数 (\(|H_{ap}(e^{j\omega})|^2 = 1\))。其零极点互为共轭倒置：
\[ H_{ap}(z) = A \prod_{k=1}^{N} \frac{z^{-1} - a_k^*}{1 - a_k z^{-1}},\quad |a_k|<1 \]
\paragraph{最小相位系统} 因果稳定，所有零极点均在单位圆内。在相同幅频响应的因果稳定系统中，其相位延迟最小。其逆系统 \(1/H(z)\) 也是因果稳定的。
\paragraph{FIR系统} \(h(n)\) 长度有限，无反馈，系统函数为 \(z^{-1}\) 的多项式；永远稳定，极点全在 \(z=0\)。
\paragraph{IIR系统} \(h(n)\) 长度无限，有反馈，系统函数为有理分数；不一定稳定。

\section{DFT}
\subsection{定义}
长度为\(N\)的序列\(x(n)\)的DFT
\[
X(k)=\sum_{n=0}^{N-1}x(n)W_N^{nk}
\]
反变换IDFT
\[
x(n)=\frac{1}{N}\sum_{k=0}^{N-1}X(k)W_N^{-nk}
\]

\subsection{性质}
\(X(k)\)周期延拓得\(\tilde{X}(k)\)，即\(\tilde{x}(n)\)的DFS

从Z变换得DFT
\[
X(k)=X(z)\bigr|_{z=W_N^k}
\]
Z域内插
\begin{gather*}
X(z)=\sum_{k=0}^{N-1}X(k)\Phi(W_N^k z),\\
\Phi(z) = \frac{1}{N}\frac{1-z^{-N}}{1-z^{-1}}
\end{gather*}

频域内插
\begin{gather*}
X(e^{i\omega})=\sum_{k=0}^{N-1}X(k)\Phi(\omega-2\pi k/N),\\
\Phi(\omega) = \frac{1}{N}\frac{\sin(N\omega/2)}{\sin(\omega/2)}e^{-j(N-1)\omega/2} = \Phi(z)\bigr|_{z=e^{j\omega}}
\end{gather*}

线性
\(
ax(n)+y(n)\Leftrightarrow aX(k)+Y(k)
\)

对偶
\(
X(n)\Leftrightarrow Nx(k)
\)

共轭
\(
x^*(n)\Leftrightarrow X^*(-k)
\)

奇偶
\(
x(-n)\Leftrightarrow X(-k)
\)

时移和频移
\begin{gather*}
x(n-n_0)\Leftrightarrow W_N^{n_0k}X(k)\\
W_N^{-nk_0}x(n)\Leftrightarrow X(k-k_0)
\end{gather*}

内积
\(
\sum_{n=0}^{N-1} x(n) y^*(n) = \frac{1}{N} \sum_{k=0}^{N-1} X(k) Y^*(k)
\)

Parseval 定理
\(
\sum_{n=0}^{N-1} |x(n)|^2 = \frac{1}{N} \sum_{k=0}^{N-1} |X(k)|^2
\)

循环卷积定义
\[
x(n)\otimes y(n)=\sum_{k=0}^{N-1}x(k)y(n-k)
\]

循环卷积定理
\begin{gather*}
x(n) \otimes y(n) \Leftrightarrow X(k) Y(k) \\
x(n) y(n) \Leftrightarrow \frac{1}{N} X(k) \otimes Y(k)
\end{gather*}

\subsection{计算线性卷积}
\paragraph{重叠相加法}
计算 \(x(n)\ast h(n)\)，\(h(n)\) 长度为 \(M\)。将 \(x(n)\) 分为长度为 \(N\) 的不重叠段，每段补零后用 \(N+M-1\) 点循环卷积计算线性卷积，结果重叠相加。
\paragraph{重叠保留法}
计算 \(x(n)\ast h(n)\)，\(h(n)\) 长度 \(M\)。将 \(x(n)\) 分为长度 \(N\) 的段，段间重叠 \(M-1\) 点。每段做 \(N\) 点循环卷积，每段输出保留后 \(N-M+1\) 点，丢弃前 \(M-1\) 点，直接拼接。

\subsection{频谱分析}
取样频率\(f_s\)，做\(N\)点DFT时

周期\(M\)信号频谱泄漏不发生当且仅当\(M|N\)

谱线\(X(k)\)代表的频率\(\omega_k=\frac{2\pi k}{N}\)

频率分辨率
\(\Delta f = f_s/N\)

最小频率偏差\(\Delta \omega=\min_{k} |\frac{2\pi}{M} - \frac{2\pi k}{N}|\)

\section{基2-FFT}
长度\(N\)是\(2\)的整数幂。所需复数乘法次数为\(N/2\log_2 N\)。
\subsection{DIT-FFT}
分成奇偶\(x_1(r)=x(2r),x_2(r)=x(2r+1)\)。
\begin{gather*}
X(k)=X_1(k)+W_N^k X_2(k)\\
X(k+N/2)=X_1(k)-W_N^k X_2(k)
\end{gather*}
\subsection{DIF-FFT}
\begin{gather*}
X(2r)=\FFT\{x(n)+x(n+N/2)\}\\
X(2r+1)=\FFT\{[x(n)-x(n+N/2)]W_N^n\}
\end{gather*}
\subsection{IFFT}
共轭法
\(
x(n)=1/N{(\FFT\{X^*(k)\})}^*
\)

\end{multicols}
\end{document}
